\documentclass[10pt,a4paper]{amsart}
\usepackage[margin=19mm]{geometry}
\usepackage{amsmath,amssymb,mathtools}
\usepackage[scr=rsfs,cal=euler]{mathalfa}
\usepackage{xfrac}
\usepackage[unicode,hidelinks,bookmarks=false]{hyperref}
\newcommand{\ie}{i.e.\ }
\newcommand{\cf}{cf.\ }
\newcommand{\resp}{respectively\ }
\newcommand{\Subsection}{\subsection}
\providecommand{\nobreakdash}{\nobreak\hbox{-}}
\setlength{\emergencystretch}{2em}
\multlinegap0pt
\usepackage{mathtools}
\usepackage{amssymb}
\usepackage[all]{xy}
\usepackage{tikz}
\usepackage{float}
\newtheorem{proposition}{Proposition}
\def\C {{\mathbb C}}
\def\R {{\mathbb R}}
\title{Hanlon-Hicks-Lazarev resolution revisited}
\author{Lev Borisov, Zengrui Han}
\date{Source: arXiv:2509.11077v2 (CC BY 4.0)}
\begin{document}
\maketitle

\noindent\emph{Source and licence.} Hanlon-Hicks-Lazarev resolution revisited. Version arXiv:2509.11077v2 (CC BY 4.0), distributed by arXiv under the Creative Commons Attribution 4.0 International licence, http://creativecommons.org/licenses/by/4.0/, as recorded in the arXiv licence metadata for this exact version and shown on the abstract page https://arxiv.org/abs/2509.11077v2. Full licence notice: \texttt{licenses/arxiv-2509.11077-CC-BY-4.0.txt}.\par\smallskip

\noindent\textit{Editorial scope.} Counted unit: the article's proposition prop-isomorphic-to-singular (source0.tex lines 238--240). The article's complete original proof of it is delivered with it (source0.tex lines 241--263, one \texttt{proof} environment of 23 lines). No mathematics is written, repaired or re-derived: the statement and the proof are the article's own lines, in the article's own notation, and no retained line is substituted or re-flowed.  Retained uncounted as the source-local context the proof needs: lines 143--143; lines 228--230.  Nothing was omitted from any retained span, so the package carries no omission record.  No retained line contains a citation, so the wrapper carries no bibliography and no bibliography entry.  No retained line contains a figure environment, an image-inclusion command or drawing code, so no image is delivered and the package declares no asset dependency.  Editorial formatting only, in addition: an AMS \texttt{amsart} preamble that restates the article's own declarations and macros used by the retained lines, namely the environments \texttt{proposition} on separate counters, together with the standard packages those lines need.  The retained environments print in this document as Proposition 1 in order of appearance, whereas the article prints the same items with the numbering of its own full document.  The complete native source of the article as distributed in the arXiv e-print for this exact version is bundled unchanged as \texttt{sources/185015/source0.tex}.
\subsection{Affine HHL complex}\label{subsec-affine-HHL}

\begin{align}\label{eq-l-complex}
    \cdots\rightarrow\bigoplus_{\substack{ \sigma\in S_m, (a_i)\in\mathbb{N}^n \\ \deg(\prod x_i^{a_i} e_{\sigma}) = l}} \C \, \prod x_i^{a_i} e_{\sigma} \xrightarrow{d} \bigoplus_{\substack{ \tau\in S_{m-1}, (a_i)\in\mathbb{N}^n\\ \deg(\prod x_i^{a_i} e_{\tau}) = l}} \C \, \prod x_i^{a_i} e_{\tau}\rightarrow\cdots.
\end{align}

\begin{proposition}\label{prop-isomorphic-to-singular}
    The $l$-component of the HHL complex is isomorphic to the cellular chain complex of $U_{\widetilde{l}}$ induced by the Bondal stratification.
\end{proposition}

\begin{proof}
    First of all, we claim that there is a natural bijection between the set of monomials $\prod x_i^{a_i} e_{\sigma}$ with degree $l$ and the set of strata $\widetilde{\sigma}$ on the universal cover $\Lambda^*_{\R}$ that are contained in $U_{\widetilde{l}}$. To see this, recall that $\widetilde{l} = \sum l_i v_i^*$ is a fixed lift of $l$ in $L^*$. Then for each monomial $\prod x_i^{a_i} e_{\sigma}$, there is a unique lift $\widetilde{\sigma}$ such that the following equality holds in $L^*$:
    \begin{align}\label{eq1}
        \sum_{i=1}^n a_i v_i^* + \sum_{i=1}^n \lceil H_i(\widetilde{\sigma}) \rceil v_i^* = \sum_{i=1}^n l_i v_i^*.
    \end{align}
    The condition $a_i\geq 0$ then ensures $\widetilde{\sigma}\subseteq U_{\widetilde{l}}$. On the other hand, for any $\widetilde{\sigma}\in \widetilde{S_m}$ with $\widetilde{\sigma}\subseteq U_{\widetilde{l}}$, it is straightforward to check that
    \begin{align*}
        \prod x_i^{l_i - \lceil H_i(\widetilde{\sigma}) \rceil } e_{\sigma}
    \end{align*}
    corresponds to $\widetilde{\sigma}$ under the correspondence above. The condition $\widetilde{\sigma}\subseteq U_{\widetilde{l}}$ ensures all exponents are non-negative. This gives maps between the complex \eqref{eq-l-complex} and
    \begin{align*}
        \cdots\rightarrow  \bigoplus_{\substack{ \widetilde{\sigma}\in \widetilde{S_m} \\ \widetilde{\sigma}\subseteq U_{\widetilde{l}} }} \C \cdot e_{\widetilde{\sigma}}\xrightarrow{\partial}  \bigoplus_{\substack{ \widetilde{\tau}\in \widetilde{S_{m-1}} \\ \widetilde{\tau}\subseteq U_{\widetilde{l}} }} \C \cdot  e_{\widetilde{\tau}}\rightarrow \cdots
    \end{align*}
    that are inverse to each other. 

    \smallskip
    
    It then suffices to check the compatibility with the differentials. Fix a monomial $\prod x_i^{a_i} e_{\sigma}$ and recall the definition of $d$ from Section \ref{subsec-affine-HHL}. Consider the set $\{\widetilde{\tau_1},\cdots,\widetilde{\tau_r}\}$ of all facets of $\widetilde{\sigma}$ whose image under the quotient map is a fixed facet $\tau$ of $\sigma$. Denote $\epsilon_{i}^{(j)}:=\left\lceil H_{i}(\widetilde{\sigma})  \right\rceil - \left\lceil H_{i}(\widetilde{\tau_j})  \right\rceil $ for all $i=1,\cdots,n$, the monomial $\prod x_i^{a_i} e_{\sigma}$ is mapped to $\left(\sum_{j=1}^r \prod x_i^{a_i+\epsilon_{i}^{(j)}}\right) e_{\tau}$. Now denote the image of $\prod x_i^{a_i+\epsilon_{i}^{(j)}}e_{\tau}$ by $\widehat{\tau_j}$. Then by definition we have
    \begin{align*}
        \sum_{i=1}^n (a_i +\epsilon_{i}^{(j)} )v_i^* + \sum_{i=1}^n \lceil H_i(\widehat{\tau_j}) \rceil v_i^* = \sum_{i=1}^n l_i v_i^*.
    \end{align*}
    A comparison with \eqref{eq1} yields $\left\lceil H_{i}(\widetilde{\tau_j})  \right\rceil = \left\lceil H_{i}(\widehat{\tau_j})  \right\rceil$ for all $i=1,\cdots,n$. This forces $\widetilde{\tau_j} = \widehat{\tau_j}$, which proves the compatibility.
\end{proof}

\end{document}
